%
% einleitung.tex -- Beispiel-File für die Einleitung
%
% (c) 2020 Prof Dr Andreas Müller, Hochschule Rapperswil
%
% !TEX root = ../../buch.tex
% !TEX encoding = UTF-8
%
\section{Einführung}
\label{basel:section:einfuehrung}
 
Was kommt heraus wenn man
\[
    1 + \frac{1}{4} + \frac{1}{9} + \frac{1}{16} + \frac{1}{25} + \cdots
\]
summiert?

Auf den ersten Blick erkennt man nicht gleich den Zusammenhang der Brüche in diesem Term.
Wenn man ihn aber umschreibt zu
\[
    \frac{1}{1^2} + \frac{1}{2^2} + \frac{1}{3^2} + \frac{1}{4^2} + \frac{1}{5^2} + \cdots
\]
erkennt man schnell das Muster.
Das Resultat dieser Reihe ist die unendliche Summe der Kehrwerte der Quadratzahlen,
\[
    \sum_{n=1}^{\infty} \frac{1}{n^2}
    =
    \frac{1}{1^2} + \frac{1}{2^2} + \frac{1}{3^2} + \frac{1}{4^2} + \frac{1}{5^2} + \cdots.
\]
Die Frage, was bei dieser Rechnung herauskommt, ist unter dem Namen \emph{Basel-Problem} bekannt.